%!TEX root = ../main.tex
\section{Experiments}
\label{sec:experiments}

We organize the evaluation around eight questions. Exp-1 asks whether \sys{} outperforms state-of-the-art memory systems across memory abilities. Exp-2 asks whether the advantage survives a change of backbone. Exp-3 asks which query patterns the advantage comes from. Exp-4 asks what the accuracy costs in LLM calls. Exp-5 asks how much each architectural layer contributes. Exp-6 asks whether the second-order memory improves a user over time. Exp-7 asks how accurate the zero-cost planner is and how sensitive the end-to-end result is to its two free parameters. Exp-8 asks how read-path cost behaves as history length grows.

\subsection{Experimental Setup}
\label{subsec:setup}

%!TEX root = ../main.tex
\begin{table}[t]
    \centering
    \small
    \caption{The eight benchmarks, grouped by the memory ability they stress.
      History length spans two orders of magnitude, which lets us separate gains
      that come from structure from gains that come from context length.}
    \label{tbl:benchmarks}
    \setlength{\tabcolsep}{4pt}
    \begin{tabular}{@{}llr@{}}
        \toprule
        \textbf{Benchmark} & \textbf{Memory ability} & \textbf{History} \\
        \midrule
        \locomo~\cite{locomo}         & Multi-session dialogue    & $\sim$9K \\
        \lmev~\cite{longmemeval}      & Entity-centric recall     & $\sim$115K \\
        \membench~\cite{membench}     & Typed access primitives   & $\sim$100K \\
        \addlinespace[2pt]
        \lmeo~\cite{longmemeval}      & Long-horizon retention    & $\sim$40K \\
        \lmes~\cite{longmemeval}      & Long-horizon retention    & $\sim$115K \\
        \lmem~\cite{longmemeval}      & Long-horizon retention    & $\sim$1.5M \\
        \addlinespace[2pt]
        \evomem~\cite{evomemory}      & Evolving user state       & $\sim$50K \\
        \dynamicmem~\cite{dynamicmem} & Evolving user state       & $\sim$127K \\
        \bottomrule
    \end{tabular}
\end{table}


\stitle{Datasets.}
We use the eight long-term memory benchmarks in Table~\ref{tbl:benchmarks}, chosen to separate three memory abilities that a single benchmark would confound. \locomo~\cite{locomo} and \lmev~\cite{longmemeval} test recall over multi-session dialogue. \lmeo, \lmes, and \lmem~\cite{longmemeval} share one query style and scale the surrounding history by nearly two orders of magnitude, which lets history length be varied as a controlled factor. \evomem~\cite{evomemory} and \dynamicmem~\cite{dynamicmem} issue many questions per user over an evolving state, the only setting in which per-user adaptation becomes observable. \membench~\cite{membench} labels every one of its $300$ queries with the query pattern it requires, and it is therefore the one benchmark on which a per-class decomposition is possible.

\stitle{Baselines.}
We compare against seven systems in two groups. \etitle{(1) Prompting Baselines.} \textbf{Full-Context}~\cite{longcontext} places the entire history in the context window, and \textbf{RAG}~\cite{rag} retrieves top-$k$ passages by dense similarity. \etitle{(2) Memory Systems.} \textbf{MemGPT}~\cite{memgpt} pages facts between context tiers. \textbf{Mem0}~\cite{mem0} consolidates facts into a dictionary through an \textsc{Add}/\textsc{Update}/\textsc{Delete} loop, and \textbf{Mem0}$^{g}$ adds entity relations. \textbf{Zep}~\cite{zep} maintains a temporal knowledge graph. \textbf{LongMemEval-Fw}~\cite{longmemeval} applies a three-stage indexing, retrieval, and reading pipeline. All seven are run from their released implementations under the same backbone, which leaves the memory design as the sole source of any difference.

\stitle{Evaluation Metrics.}
Multiple-choice benchmarks are scored by exact match. The remaining benchmarks are scored by F1 against the gold answer, normalized at the phrase level for \lmev{} and at the token level elsewhere. Phrase-level normalization credits a prediction only when it reproduces the gold phrase, so absolute F1 on \lmev{} sits far below the token-level scores, and the two scales are not comparable; within \lmev{} the ranking across systems is unaffected. Higher is better throughout, and every reported number is an accuracy on the benchmark's own scale.

\stitle{Backbone LLMs and Settings.}
Unless stated otherwise all systems run on GPT-5.5. Exp-2 additionally evaluates three further closed-source backbones~\cite{gpt4,claude,gemini} and four open-source backbones, selected to span a wide capability range at two deployment tiers. \sys{} is implemented in Python against an OpenAI-compatible client. Every hyperparameter is fixed once and reused across all benchmarks and backbones, with no per-dataset tuning.

%!TEX root = ../main.tex
\begin{table*}[t]
    \centering
    \small
    \caption{Accuracy across eight long-term memory benchmarks. \sys{} matches
      or improves the strongest prior system on every benchmark, and the margin
      is widest on the long-horizon and state-inference groups where a single
      retrieval pass is least sufficient.}
    \label{tbl:main_results}
    \begin{tabular}{@{}l ccc ccc cc@{}}
        \toprule
        \multirow{2}{*}{\textbf{Method}}
            & \multicolumn{3}{c}{\textbf{Conversational memory}}
            & \multicolumn{3}{c}{\textbf{Long-horizon memory}}
            & \multicolumn{2}{c}{\textbf{Personalized state}} \\
        \cmidrule(lr){2-4} \cmidrule(lr){5-7} \cmidrule(l){8-9}
            & \textbf{\locomo} & \textbf{\lmev} & \textbf{\membench}
            & \textbf{\lmeo} & \textbf{\lmes} & \textbf{\lmem}
            & \textbf{\evomem} & \textbf{\dynamicmem} \\
        \midrule
        \tblgroup{9}{Prompting baselines} \\
        Full-Context~\cite{longcontext}   & 52.3 & 4.5 & 86.7 & 49.3 & 43.6 & 58.5 & 16.7 & 33.5 \\
        RAG~\cite{rag}                    & 53.1 & 4.9 & 87.8 & 50.8 & 45.1 & 59.1 & 17.3 & 34.1 \\
        \addlinespace[2pt]
        \tblgroup{9}{Memory systems} \\
        MemGPT~\cite{memgpt}              & 54.5 & 6.3 & 89.7 & 54.1 & 48.4 & 61.1 & 18.0 & 34.3 \\
        Mem0~\cite{mem0}                  & 56.4 & 7.2 & \second{91.2} & 58.3 & 54.2 & 65.2 & 20.4 & 34.5 \\
        Mem0$^{g}$~\cite{mem0}            & 57.1 & 7.5 & 90.8 & 60.1 & 56.7 & 67.1 & 21.2 & 34.7 \\
        Zep~\cite{zep}                    & \second{58.2} & 8.1 & 90.4 & \second{62.2} & \second{59.3} & \second{70.3} & \second{22.1} & \second{34.9} \\
        LongMemEval-Fw~\cite{longmemeval} & 57.3 & \second{9.4} & 90.1 & 61.7 & 58.8 & 69.4 & 21.6 & 34.6 \\
        \addlinespace[2pt]
        \ourgroup{9}{Ours} \\
        \textbf{\sys{}}
            & \best{58.8} & \best{10.3} & \best{92.8}
            & \best{67.2} & \best{64.9} & \best{75.2}
            & \best{24.6} & \best{34.9} \\
        \bottomrule
    \end{tabular}
\end{table*}


\subsection{Exp-1. Accuracy Across Three Memory Abilities}
\label{subsec:exp1}

We compare \sys{} against all seven baselines on the eight benchmarks of Table~\ref{tbl:benchmarks}, grouped by the memory ability each one stresses. Table~\ref{tbl:main_results} gives the comparison. \sys{} matches or improves the strongest prior system on every benchmark, and the margin varies systematically with the memory ability the benchmark stresses.

\stitle{Long-Horizon Memory.}
The margin is widest on the three LongMemEval variants, where it reaches $4.9$ to $5.6$ points, and it does not shrink as history grows by nearly two orders of magnitude. The reason lies in what each system hands the model. A longer history enlarges the candidate set that similarity retrieval returns, so prior systems pass the model more passages and leave it to locate, order, and align the relevant facts inside the context window, the regime in which long-context models are known to undercount and misorder~\cite{longcontext,ruler}. \sys{} reads the fields a query needs from the catalog indexes, so the context block stays small and already structured, and the model only performs the reasoning step. The size of that block is set by the query pattern, independent of $|\mathcal{H}|$, which is why the margin persists at the largest scale.

\stitle{Conversational Memory.}
The margin is narrowest on this group, at $0.6$ points on \locomo{} and $1.6$ on \membench{}. These benchmarks ask mostly for facts stated in one place, a query pattern that a single retrieval pass already serves well, so there is less headroom for an operator to recover. The result on \lmev{} is $10.3$ against $9.4$ for the strongest baseline, a $9.6\%$ relative improvement on a strict phrase-level F1 scale on which every system scores in the single digits. What distinguishes \sys{} on this group is that the accuracy and the cost advantage hold together: on \membench{} it reaches the highest accuracy while issuing fewer LLM calls than the uniform reference, as Exp-4 shows.

\stitle{Personalized State.}
The two state-inference benchmarks separate the systems least in absolute terms, with all seven baselines falling within $5.4$ points on \evomem{} and $1.4$ points on \dynamicmem. These benchmarks issue many questions per user, and a system that treats each query independently answers the tenth exactly as it answered the first, which is what prior designs do: they record what arrives on the write path and never what succeeded on the read path. An aggregate accuracy therefore averages over query position, the dimension on which the second-order memory acts, and hides any per-user improvement. Exp-6 conditions on query position and recovers the effect that the aggregate hides.

\finding{\sys{} matches or improves the strongest prior system on all eight benchmarks, and the margin is largest where a single retrieval pass is least sufficient, on long histories and on repeated questions from one user.}

\subsection{Exp-2. Backbone Generalization}
\label{subsec:exp2}

A system that supplies external structure should benefit every backbone, and the benefit should be largest where the backbone's own capacity is most limited. To test this, we hold \membench{} and the full \sys{} configuration fixed and vary only the backbone, over four closed-source and four open-source models spanning a wide capability range. Table~\ref{tbl:main_results_open} collects the comparison.

%!TEX root = ../main.tex
\begin{table}[t]
    \centering
    \small
    \caption{Accuracy on \membench{} across backbones. \textbf{Best prior} is
      the strongest of the seven baselines of Table~\ref{tbl:main_results}
      re-run on that same backbone. The gain is smallest where the backbone is
      strongest, so the structure \sys{} supplies substitutes for capacity the
      model does not have.}
    \label{tbl:main_results_open}
    \begin{tabular}{@{}lcc@{}}
        \toprule
        \textbf{Backbone} & \textbf{Best prior} & \textbf{\sys{}} \\
        \midrule
        \tblgroup{3}{Closed-source} \\
        GPT-5.5             & 91.2 & \best{92.8} \\
        Claude-3.5-Sonnet   & 87.4 & \best{91.4} \\
        Gemini-1.5-Pro      & 85.6 & \best{89.7} \\
        GPT-4o              & 82.8 & \best{88.7} \\
        \addlinespace[2pt]
        \tblgroup{3}{Open-source} \\
        DeepSeek-V4         & 82.7 & \best{88.5} \\
        Llama-3.1-70B       & 79.8 & \best{86.5} \\
        Qwen3-14B           & 77.4 & \best{85.1} \\
        Qwen3-8B            & 73.4 & \best{81.3} \\
        \bottomrule
    \end{tabular}
\end{table}


\sys{} improves every backbone, which rules out an explanation resting on one model's inductive bias. The size of the improvement is what varies. It is smallest on the strongest backbone and largest on the weakest, and this ordering holds within each group separately. This pattern suggests that external structure compensates for capability the backbone lacks. A frontier model already holds more history in context and counts over it more reliably, so the incremental gain is smaller, but the gain remains positive even on the strongest backbone. The practical consequence is that \sys{} widens the set of backbones on which long-term memory is usable, which matters most where frontier inference per query is not affordable.

\finding{The advantage of \sys{} is not backbone-specific, and it grows as backbone capability falls, with external memory structure substituting for internal model capacity.}

\subsection{Exp-3. Per-Pattern Accuracy Decomposition}
\label{subsec:exp3}

%!TEX root = ../main.tex
\begin{table}[t]
    \centering
    \small
    \caption{Accuracy by plan class on \membench. Every system handles point
      lookups well, and the systems separate on aggregation and comparison,
      which is where an explicit plan replaces free-form reasoning.}
    \label{tbl:perclass}
    \begin{tabular}{@{}lccc@{}}
        \toprule
        \textbf{Method} & \opPOINT & \opAGG & \opCMP \\
        \midrule
        \tblgroup{4}{Prompting baselines} \\
        Full-Context~\cite{longcontext}    & 90.9 & 64.4 & 84.2 \\
        RAG~\cite{rag}                     & 91.8 & 67.1 & 85.8 \\
        \addlinespace[2pt]
        \tblgroup{4}{Memory systems} \\
        MemGPT~\cite{memgpt}               & 92.9 & 68.5 & 87.0 \\
        Mem0~\cite{mem0}                   & 93.9 & 70.2 & \second{88.1} \\
        Mem0$^{g}$~\cite{mem0}             & 94.1 & \second{71.9} & 85.4 \\
        Zep~\cite{zep}                     & \second{94.3} & 71.2 & 85.4 \\
        LongMemEval-Fw~\cite{longmemeval}  & 93.8 & 70.8 & 84.7 \\
        \addlinespace[2pt]
        \ourgroup{4}{Ours} \\
        \textbf{\sys{}} & \best{94.6} & \best{78.0} & \best{91.2} \\
        \bottomrule
    \end{tabular}
\end{table}


Exp-1 established that \sys{} leads overall. To locate the source of that lead, we decompose \membench{} accuracy by the pattern each query requires, which \membench{} labels directly. Table~\ref{tbl:perclass} breaks the accuracy out by query pattern.

The decomposition reveals where the advantage originates. \sys{} leads on all three query patterns, and the gains concentrate on aggregation and comparison, where a specialized operator replaces unstructured reasoning. The separation appears on these two classes because they require the model to compose evidence across multiple records.

\stitle{Aggregation.}
Counting fails for reasons that a better retriever cannot fix. When candidate items are scattered through a long context, models miscount even after retrieving every relevant item~\cite{ruler}. Prior systems retrieve passages and ask for a count, which leaves the count implicit inside unstructured generation. \sys{} makes it explicit, enumerating candidates, deciding each one against the predicate, then tallying. Each step is bounded, and the total rests on enumerable evidence with an exact count.

\stitle{Comparison.}
Comparison fails through misalignment. The model must locate the compared attribute in two separate passages and align them before it can order them, and that alignment is what breaks on long histories. \sys{} performs the alignment outside the model, reading the attribute from the typed index for each candidate and placing both values adjacent in the prompt. The remaining task is an ordering over two values.

Both failure modes live in how the model \emph{processes} retrieved evidence, which is why improving retrieval alone does not close the gap and why the computation has to follow from the query pattern.

%!TEX root = ../main.tex
\begin{table}[t]
    \centering
    \small
    \caption{One aggregation query traced through three systems. The two
      baselines fail differently, and each failure is removed by a different
      layer of \sys{}.}
    \label{tbl:casestudy}
    \begin{tabular}{@{}llcc@{}}
        \toprule
        \textbf{System} & \textbf{Retrieval and reasoning} & \textbf{Calls} & \textbf{Answer} \\
        \midrule
        Full-Context   & count left implicit & 1 & \textcolor{dropred}{5} \\
        Zep~\cite{zep} & predicate unresolved & 1 & \textcolor{dropred}{4} \\
        \midrule
        \ourgroup{4}{\sys{}, layer by layer} \\
        \multicolumn{4}{@{}l}{\quad \emph{Planner} routes on the quantifier marker} \\
        \multicolumn{4}{@{}l}{\quad \emph{Second-order memory} fronts this user's resolved entities} \\
        \multicolumn{4}{@{}l}{\quad \emph{Operator} decides each candidate, then tallies} \\
        \multicolumn{4}{@{}l}{\quad \emph{Executor} withholds the third call on agreement} \\
        \textbf{\sys{}} & explicit per-item decision & 2 & \textcolor{cadmiumgreen}{\textbf{3}} \\
        \bottomrule
    \end{tabular}
\end{table}


To make the two failure modes concrete, Table~\ref{tbl:casestudy} traces one aggregation query through \sys{} and two baselines. Both baselines had the relevant material available, and in neither case did retrieval fail. The prompting baseline miscounts evidence spread across sessions, and the graph memory retrieves topically related nodes without resolving the predicate the query asks about. \sys{} removes the two failures with different layers. The planner selects the \agg{} pattern through Eq.~(\ref{eq:classifier}), the second-order memory promotes the entities this user's earlier queries resolved, the operator records an explicit decision per candidate before tallying, and the executor withholds its third call once two samples agree.

\finding{The accuracy advantage concentrates on aggregation and comparison, where a specialized operator replaces unstructured reasoning over scattered evidence.}

%!TEX root = ../main.tex
\begin{table}[t]
    \centering
    \small
    \caption{Calls per query on \membench. The two \sys{} rows share an accuracy
      column, which isolates the short-circuit as a pure cost mechanism and
      attributes the accuracy to the plan and not to the budget.}
    \label{tbl:cost}
    \setlength{\tabcolsep}{3pt}
    \begin{tabular}{@{}lcccrc@{}}
        \toprule
        \multirow{2}{*}{\textbf{Configuration}}
            & \multicolumn{3}{c}{\textbf{Calls per query}}
            & \multirow{2}{*}{\textbf{Total}} & \multirow{2}{*}{\textbf{Acc.}} \\
        \cmidrule(lr){2-4}
            & \opPOINT & \opAGG & \opCMP & & \\
        \midrule
        Uniform, one call            & 1.00 & 1.00 & 1.00 & 300 & 86.7 \\
        Uniform, three calls         & 3.00 & 3.00 & 3.00 & 900 & 89.1 \\
        \addlinespace[2pt]
        \sys{}, budget only          & 1.00 & 3.00 & 3.00 & 532 & 92.8 \\
        \textbf{\sys{}}              & 1.00 & 2.30 & 2.18 & \best{444} & \best{92.8} \\
        \bottomrule
    \end{tabular}
\end{table}

%!TEX root = ../../main.tex
% Cost-accuracy Pareto plot on MemBench.
\begin{figure}[t]
    \centering
    \resizebox{0.98\columnwidth}{!}{%
    \begin{tikzpicture}[
        every node/.style={font=\footnotesize},
        axis/.style={-Latex, thick, black},
        tick/.style={black, thin},
        grid/.style={gray!25, dashed, very thin},
        pt/.style={circle, minimum size=6pt, inner sep=0pt},
        lead/.style={gray!55, thin},
    ]
        \def\W{6.4}
        \def\H{3.6}
        %-------- axes with arrowheads only --------%
        \draw[axis] (0,0) -- (\W+0.15,0);
        \draw[axis] (0,0) -- (0,\H+0.15);
        %-------- x ticks --------%
        \foreach \x/\lbl in {0.0/1.0, 1.4/1.5, 2.8/2.0, 4.2/2.5, 5.6/3.0}{
            \draw[tick] (\x,0) -- (\x,-0.09);
            \node[below] at (\x,-0.09) {\lbl};
        }
        %-------- y ticks --------%
        \foreach \y/\lbl in {0.0/86, 0.9/88, 1.8/90, 2.7/92, 3.6/94}{
            \draw[tick] (0,\y) -- (-0.09,\y);
            \node[left] at (-0.09,\y) {\lbl};
        }
        %-------- axis labels --------%
        \node[below, font=\footnotesize] at (\W/2,-0.7) {\emph{Average LLM calls per query}};
        \node[rotate=90, font=\footnotesize] at (-0.9,\H/2) {\emph{Accuracy on \membench{} (\%)}};
        %-------- soft grid --------%
        \foreach \x in {1.4, 2.8, 4.2, 5.6} \draw[grid] (\x,0) -- (\x,\H);
        \foreach \y in {0.9, 1.8, 2.7, 3.6} \draw[grid] (0,\y) -- (\W,\y);
        %-------- data points --------%
        % x = (calls-1.0)*2.8;  y = (acc-86)*0.45
        \node[pt, fill=gray!70]    at ({(1.00-1.0)*2.8}, {(86.7-86)*0.45}) (fc)   {};
        \node[pt, fill=gray!70]    at ({(1.00-1.0)*2.8}, {(87.8-86)*0.45}) (rag)  {};
        \node[pt, fill=softblue!70] at ({(1.4-1.0)*2.8}, {(89.7-86)*0.45}) (mgpt) {};
        \node[pt, fill=softblue!70] at ({(1.5-1.0)*2.8}, {(91.2-86)*0.45}) (m0)   {};
        \node[pt, fill=softblue!70] at ({(1.8-1.0)*2.8}, {(90.8-86)*0.45}) (m0g)  {};
        \node[pt, fill=softblue!70] at ({(1.7-1.0)*2.8}, {(90.4-86)*0.45}) (zep)  {};
        \node[pt, fill=softblue!70] at ({(1.6-1.0)*2.8}, {(90.1-86)*0.45}) (lfw)  {};
        \node[pt, fill=softorange] at ({(3.00-1.0)*2.8}, {(89.1-86)*0.45}) (uni)  {};
        \node[pt, fill=dropred, minimum size=8pt] at ({(1.48-1.0)*2.8},{(92.8-86)*0.45}) (mc) {};
        %-------- pareto arrow --------%
        \draw[dropred, dashed, thick, opacity=0.5] (fc.center) -- (mc.center);
        %-------- labels with leaders, dispersed --------%
        \node[anchor=west] (mcL) at (2.3, 3.35) {\textbf{\sys{} (ours)}};
        \draw[lead] (mc) -- (mcL.west);
        \node[anchor=west] (m0L)  at (2.5, 2.60) {Mem0 (91.2)};
        \draw[lead] (m0) -- (m0L.west);
        \node[anchor=west] (m0gL) at (2.9, 2.15) {Mem0$^{g}$ (90.8)};
        \draw[lead] (m0g) -- (m0gL.west);
        \node[anchor=west] (zepL) at (2.6, 1.80) {Zep (90.4)};
        \draw[lead] (zep) -- (zepL.west);
        \node[anchor=west] (lfwL) at (2.3, 1.45) {LongMemEval-Fw (90.1)};
        \draw[lead] (lfw) -- (lfwL.west);
        \node[anchor=west] (mgptL) at (2.0, 1.10) {MemGPT (89.7)};
        \draw[lead] (mgpt) -- (mgptL.west);
        \node[anchor=east] (uniL) at (5.5, 0.60) {Uniform-$K{=}3$ (89.1)};
        \draw[lead] (uni) -- (uniL.east);
        \node[anchor=west] (ragL) at (0.25, 0.90) {RAG (87.8)};
        \draw[lead] (rag) -- (ragL.west);
        \node[anchor=west] (fcL) at (0.25, 0.30) {Full-Context (86.7)};
        \draw[lead] (fc) -- (fcL.west);
    \end{tikzpicture}
    }
    \caption{Cost--accuracy Pareto on \membench{} (GPT-5.5, $n{=}300$). Each point is a system's average LLM calls per query (x-axis) versus its accuracy (y-axis). \sys{} occupies the upper-left corner: strictly higher accuracy than every prior memory system and strictly lower cost than the Uniform-$K{=}3$ reference. No prior system Pareto-dominates \sys{}.}
    \Description{Scatter plot of average LLM calls per query against accuracy on
      MemBench. MemCatalog sits in the upper-left region, with higher accuracy
      than every prior memory system and lower cost than the uniform K equals
      three reference.}
    \label{fig:pareto}
\end{figure}


\subsection{Exp-4. Cost Analysis}
\label{subsec:exp4}

Accuracy bought with extra LLM calls is an artifact of the budget, and the architecture should be evaluated separately. To separate the two, we measure the LLM calls each system issues per query on \membench{} and place cost and accuracy on the same axes. Cost is counted on the read path for every system, so a system that spends a retrieval or rewriting call before generating is charged for it. Write-path extraction is a one-time per-session cost that every structured system amortizes over the queries touching that session, including the on-demand extraction of \sys{}, and it is excluded uniformly for all systems. We additionally include two reference configurations that fix the sampling budget uniformly at one call and at three calls per query, which isolates the effect of \emph{allocating} a budget from the effect of its \emph{size}.

Figure~\ref{fig:pareto} gives the result. \sys{} occupies the upper-left region, where no evaluated system attains both lower cost and comparable accuracy. The two groups miss it for different reasons. Prior memory systems are inexpensive and lose accuracy on the queries a single pass cannot answer. The uniform three-call reference recovers part of that accuracy and remains dominated, because it spends its extra samples on the point lookups that never required them.

Table~\ref{tbl:cost} separates the two mechanisms behind this position. Pattern-aware budgeting acts across queries, assigning the single-call budget to the point lookups that form the majority of the workload. Agreement-based early stopping, the rule of Eq.~(\ref{eq:short}), acts within the remaining queries and stops early on those whose two samples already agree. The first mechanism draws from the easy pattern and the second from the easy instances of the hard patterns, and the slices are disjoint, which lets the savings compose. The accuracy column settles what the savings cost: it is unchanged between the two \sys{} configurations, so early stopping removes $88$ of $532$ calls without giving up a single answer, and the accuracy of \sys{} is attributable to the framework.

\finding{\sys{} is Pareto-optimal against every evaluated system, and early stopping lowers its cost at unchanged accuracy, which attributes the accuracy to the framework, with early stopping purely reducing cost.}

%!TEX root = ../main.tex
\begin{table}[t]
    \centering
    \small
    \caption{Removing one component at a time. Each layer contributes on a
      distinct axis, accuracy for the framework, accuracy under repeated
      queries for the second-order memory, and cost for the executor,
      confirming the three are complementary.}
    \label{tbl:ablation}
    \begin{tabular}{@{}lccccc@{}}
        \toprule
        \multirow{2}{*}{\textbf{Configuration}}
            & \multicolumn{2}{c}{\textbf{\membench}}
            & \multicolumn{2}{c}{\textbf{\lmev}} \\
        \cmidrule(lr){2-3} \cmidrule(l){4-5}
            & \textbf{Acc.} & \textbf{Calls} & \textbf{Acc.} & \textbf{Calls} \\
        \midrule
        \textbf{\sys{}}               & \best{92.8} & \best{444} & \best{10.3} & \best{602} \\
        \addlinespace[2pt]
        \tblgroup{5}{Framework} \\
        w/o planner          & 89.1 & 900 & 8.8 & 903 \\
        w/o physical operators        & 87.9 & 900 & 8.3 & 903 \\
        w/o catalog                   & 86.7 & 900 & 7.6 & 903 \\
        \addlinespace[2pt]
        \tblgroup{5}{Second-order memory} \\
        w/o workload log              & 92.1 & 444 & 8.6 & 602 \\
        w/o index warming             & 92.5 & 444 & 9.3 & 602 \\
        w/o session eviction          & 92.8 & 444 & 10.1 & 651 \\
        \addlinespace[2pt]
        \tblgroup{5}{Cost-aware executor} \\
        w/o plan-class budget         & 92.5 & 756 & 10.3 & 903 \\
        w/o stability short-circuit   & 92.8 & 532 & 10.3 & 723 \\
        \bottomrule
    \end{tabular}
\end{table}


\subsection{Exp-5. Ablation Study}
\label{subsec:exp5}

The three layers of \sys{} were motivated by three different obstacles. This experiment confirms that each layer contributes independently by removing one component at a time and measuring both accuracy and cost on \membench{} and \lmev{}, a multiple-choice benchmark and an open-answer one. Table~\ref{tbl:ablation} collects the ablations.

Each layer contributes on a distinct axis, confirming they are complementary. Removing any framework component costs accuracy and simultaneously removes the pattern that the budget depends on, so cost rises to the uniform setting. Within the framework the ordering is informative. Removing the catalog is worse than removing the planner, since a planner that routes correctly still has nothing typed to read. The second-order memory contributes most on \lmev{}, where a user's repeated questions benefit from read-path feedback, and its three components are not equally active on both benchmarks: session eviction leaves \membench{} untouched, whose histories stay within the cache, while on \lmev{} it saves $49$ calls. In the executor the two mechanisms differ in kind. Pattern-aware budgeting lowers cost on both benchmarks and also contributes $0.3$ points on \membench, since a point lookup that is sampled three times can be outvoted on the one answer the index already supplies. Agreement-based early stopping changes no accuracy figure on either benchmark and is a cost mechanism alone.

\finding{Each layer contributes on a distinct axis, accuracy for the framework, accuracy under repeated queries for the second-order memory, and cost for the executor, confirming they are complementary.}

%!TEX root = ../main.tex
\begin{table}[t]
    \centering
    \small
    \caption{Accuracy against the position of a query in a user's own sequence,
      evaluated on users with $\geq 10$ queries so that all four buckets are
      populated. Only \sys{} improves as a user keeps asking. On \dynamicmem{}
      it starts below the static rankers in the first bucket, where its
      workload log is still empty, and overtakes them once the log carries
      signal. Absolute levels differ from Table~\ref{tbl:main_results}, which
      averages over all users.}
    \label{tbl:second_order_learning}
    \begin{tabular}{@{}lcccc@{}}
        \toprule
        \multirow{2}{*}{\textbf{Method}}
            & \multicolumn{4}{c}{\textbf{Query position for one user}} \\
        \cmidrule(l){2-5}
            & 1--3 & 4--6 & 7--9 & 10+ \\
        \midrule
        \tblgroup{5}{\evomem} \\
        Full-Context                   & 16.7 & 17.0 & 16.8 & 16.5 \\
        Mem0~\cite{mem0}               & 20.1 & 20.5 & 20.4 & 20.2 \\
        Zep~\cite{zep}                 & \second{21.8} & \second{22.2} & \second{22.1} & \second{22.0} \\
        \textbf{\sys{}}                & \best{22.2} & \best{24.4} & \best{26.1} & \best{27.8} \\
        \addlinespace[2pt]
        \tblgroup{5}{\dynamicmem} \\
        Full-Context                   & 28.7 & 28.8 & 28.6 & 28.4 \\
        Mem0~\cite{mem0}               & \second{29.4} & 29.8 & 29.6 & 29.3 \\
        Zep~\cite{zep}                 & \best{29.9} & \second{30.0} & \second{29.9} & \second{29.8} \\
        \textbf{\sys{}}                & 29.1 & \best{30.3} & \best{31.7} & \best{32.8} \\
        \bottomrule
    \end{tabular}
\end{table}


\subsection{Exp-6. Accuracy by Query Position}
\label{subsec:exp6}

The second-order memory predicts a specific and falsifiable behavior. Accuracy on one user should rise as that user asks more questions, and no system without a read-path feedback loop should show the same rise. To test this we restrict to users with at least ten queries, so that all four position buckets are populated, condition on the position of a query within its own user's sequence, and report accuracy per position bucket, holding backbone, prompt, and benchmark fixed so that position is the only variable. Table~\ref{tbl:second_order_learning} bins the accuracy by position. Because the evaluation population is a subset of the full test set, the absolute accuracy levels in each column differ from Table~\ref{tbl:main_results}; the within-table ranking and the monotone trend are the subject of this experiment.

The predicted separation appears on both benchmarks. \sys{} improves monotonically across position buckets, gaining $5.6$ points on \evomem{} and $3.7$ on \dynamicmem{} from the first bucket to the last, while no baseline moves by more than $0.5$ points in either direction and none shows a monotone trend, since none of them records which stored entries produced accepted answers and none has a signal that could vary with position. The starting point is equally informative. In the first bucket the workload log is still empty, so \sys{} has no second-order signal yet and begins level with the static rankers, below Zep by $0.8$ points on \dynamicmem. It overtakes them from the second bucket onward and leads by $3.0$ points by the last. A mechanism that learns from answered queries should hold no advantage before any query has been answered, and this is the crossover that behavior produces.

The improvement cannot be attributed to the model or the prompt, both of which are identical across buckets. It follows the priority score of Eq.~(\ref{eq:pi}), which accumulates evidence that an entity participates in answers this user accepts. By the later buckets the context block handed to the operator is composed largely of entities with a demonstrated history of being queried, where a static ranker still offers entities selected by similarity to the current query alone. This is the read-path signal of Section~\ref{sec:secondorder}, and Exp-5 confirms that removing the workload log removes the effect.

\finding{Only \sys{} improves as a user continues asking, which isolates read-path feedback as the mechanism and shows that prior systems are structurally unable to exhibit this behavior.}

\subsection{Exp-7. Planner Accuracy and Parameter Stability}
\label{subsec:exp7routing}

%!TEX root = ../main.tex
\begin{table}[t]
    \centering
    \small
    \caption{Query-classification quality and parameter stability on \membench. The
      upper block compares the rule-based planner against an LLM planner
      that spends one call per query and against a fixed planner. The lower
      block varies the two parameters \sys{} fixes system-wide around their
      defaults $\theta{=}0.6$ and $\alpha{=}0.5$, holding all else equal.}
    \label{tbl:routing}
    \setlength{\tabcolsep}{4pt}
    \begin{tabular}{@{}lccc@{}}
        \toprule
        \textbf{Planner} & \textbf{Class acc.} & \textbf{Planner calls} & \textbf{End-to-end acc.} \\
        \midrule
        Rule-based (ours) & 91.4 & 0 & 92.8 \\
        LLM planner          & 95.1 & 1 & 93.2 \\
        Always \point        & 61.3 & 0 & 87.5 \\
        \midrule
        \tblgroup{4}{Parameter perturbation} \\
        $\theta \in [0.4, 0.8]$  & \multicolumn{3}{c}{end-to-end acc.\ 92.3--92.9} \\
        $\alpha \in [0.3, 0.7]$  & \multicolumn{3}{c}{end-to-end acc.\ 92.2--92.8} \\
        \bottomrule
    \end{tabular}
    \Description{Planner comparison and parameter sensitivity on MemBench.}
\end{table}


The planner governs which operator handles each query, so its quality directly affects the end-to-end result. This experiment validates its design along two dimensions, classification quality and parameter stability. We measure classification accuracy against the class labels \membench{} supplies, comparing the rule-based planner against an LLM planner that spends one call per query and against a fixed planner that assigns every query to \point. We then perturb the early-stopping threshold and the leading weight of the priority score over wide intervals and re-measure end-to-end accuracy. Table~\ref{tbl:routing} covers both.

Three observations follow from the table. First, the rule-based planner reaches $91.4\%$ classification accuracy at zero LLM cost. Second, classification quality has already saturated with respect to what the system does with it: an LLM planner classifies better, at $95.1\%$, yet buys only $+0.4$ end-to-end accuracy for one extra call on every query, which is $300$ additional calls against a total of $444$, and a gap of $0.4$ is no larger than the spread the parameter perturbations below produce. Paying two-thirds more calls for a difference of that size is not a trade this workload rewards. Third, the fixed planner locates where the value of query classification actually sits: routing every query to \point{} costs $5.3$ points end-to-end, an order of magnitude more than separates any two planner variants, so what matters is that the query pattern is chosen at all, and the choice granularity has diminishing returns. The lower block perturbs the two parameters \sys{} fixes system-wide, the early-stopping threshold $\theta$ and the leading weight $\alpha$ of the priority score, across intervals spanning half of each parameter's usable range. End-to-end accuracy stays within $0.6$ points over both, which leaves the reported result independent of any parameter search.

\finding{The rule-based planner reaches $91.4\%$ classification accuracy at zero cost, an LLM planner cannot convert better classification into materially better answers, and end-to-end accuracy is stable across wide parameter intervals.}

\subsection{Exp-8. Scalability Analysis}
\label{subsec:exp8}

%!TEX root = ../main.tex
\begin{table}[t]
    \centering
    \small
    \caption{Behavior as history grows on the three LongMemEval variants.
      Read-path cost is governed by the plan class and not by history length,
      so it stays flat while the prompting baseline grows with the context
      it must carry.}
    \label{tbl:scaling}
    \setlength{\tabcolsep}{4pt}
    \begin{tabular}{@{}lccc@{}}
        \toprule
        \textbf{History} & \textbf{Calls / query} & \textbf{Latency (s)} & \textbf{Accuracy} \\
        \midrule
        \tblgroup{4}{Full-Context} \\
        $\sim$40K  (\lmeo)  & 1.0 & 2.3  & 49.3 \\
        $\sim$115K (\lmes)  & 1.0 & 5.8  & 43.6 \\
        $\sim$1.5M (\lmem)  & 1.0 & 47.6 & 58.5 \\
        \addlinespace[2pt]
        \tblgroup{4}{\sys{}} \\
        $\sim$40K  (\lmeo)  & 1.9 & 3.2 & 67.2 \\
        $\sim$115K (\lmes)  & 1.9 & 3.4 & 64.9 \\
        $\sim$1.5M (\lmem)  & 1.9 & 3.5 & 75.2 \\
        \bottomrule
    \end{tabular}
    \Description{Cost and accuracy as history length grows from 40K to 1.5M tokens.}
\end{table}


A memory system is only interesting if its cost survives the growth that motivates it. The three LongMemEval variants share one query style and scale the surrounding history by nearly two orders of magnitude, from $40$K to $1.5$M tokens, which isolates history length as the varied factor. We measure calls per query, end-to-end latency, and accuracy at each scale against the prompting baseline, whose prompt necessarily carries the history it searches. Table~\ref{tbl:scaling} tracks all three.

Cost separates the two systems along the length axis. Calls per query are flat for both, at $1.9$ for \sys{} and $1.0$ for Full-Context, so latency is where the difference appears. Latency grows $20.7\times$ for Full-Context, from $2.3$s to $47.6$s, and $1.09\times$ for \sys{}, from $3.2$s to $3.5$s. At the shortest history \sys{} is the slower of the two, by $0.9$s, because catalog access is a fixed overhead that a single prompting call does not pay. That overhead does not grow, so it is already repaid at $115$K, where \sys{} answers in $3.4$s against $5.8$s, and at $1.5$M it is $13.6\times$ faster. What the table shows is a different asymptote, with \sys{} converging to a flat cost while the baseline grows linearly. Accuracy moves with the difficulty of each variant for both systems, and the ranking does not: \sys{} leads at all three scales, by $17.9$, $21.3$, and $16.7$ points, with no erosion as history grows.

The asymmetry is structural and follows from where each system spends. A prompting baseline must carry the history it searches, so its prompt grows with $|\mathcal{H}|$ and its latency grows with it. The read-path cost of \sys{} is set by the query pattern through Eq.~(\ref{eq:cost}), which depends on the shape of the query alone, independent of the size of the store, and adding sessions changes what the catalog contains without changing what a query costs to answer.

\finding{Read-path cost is governed by the query pattern, independent of history length, which turns the growth of the store from a linear cost into a constant one and leaves the accuracy advantage of \sys{} intact at every scale.}
